\section{CONCLUDING REMARKS}

\begin{onehalfspace}
This thesis examines the Turkey's labor market dynamics over the past
two decades, leveraging the comprehensive micro data provided by TurkStat.
The primary focus of this study is to unravel and analyze the intricate
interplay surrounding the phenomenon of job polarization within the
Turkish context. Through the utilization of both descriptive and empirical
analyses, a multitude of noteworthy observations come to the fore,
demanding comprehensive examination.

Firstly, a pivotal observation emerges with the stagnation in real
wage growth for the high-skilled segment of the Turkish workforce
after the year 2016. This intriguing pattern diverges from the conventional
trajectory, wherein technological advancements typically augment the
skill premium. The necessity to investigate to what extent the pattern
in Turkey's labor market is attributed to the minimum wage hikes becomes
evident.

Secondly, a distinct polarization pattern becomes evident during the
latter decade of the analyzed dataset, encompassing the most recent
ten years. This pattern manifests as employment growth at the tail
of the wage distribution and a systematic decline in employment within
the middle range. This structural phenomenon warrants particular attention
within the context of this study, prompting a comprehensive exploration
to uncover the underlying drivers of this polarization pattern.

In pursuit of explanatory insights, firstly, this study employs a
shift-share analysis. This analytical approach is wielded to dissect
the shifts in employment dynamics across various sectors. One prominent
insight derived from this examination is the role of the transition
from the manufacturing sector to the service sector. This transition
contributes partially to the observed polarization pattern. Specifically,
as employment migrates from manufacturing to the service sector, it
appears to be channeled into both the low-skilled service segment,
characterized by inherent pressures, as well as the high-skilled segment. 

In conclusion, this investigation provides valuable insights into
Turkey's labor market dynamics, particularly within the context of
job polarization. The observed wage stagnation among highly skilled
workers and the polarized employment growth patterns necessitate a
nuanced understanding of the intricate forces at play. By shedding
light on the interplay between sectoral shifts, employment dynamics,
and wage distribution, this study contributes to the broader discourse
surrounding labor market transformations and contributes with the
evidence from Turkey.
\end{onehalfspace}
\begin{center}
\pagebreak
\par\end{center}
